\chapter{Conclusion}
\label{ch:conc}

The need for theoretical predictions within QCD has a significance importance to
help connect us to experiment. Throughout the course of this thesis, we have 
demonstrated exactly this through two research projects which outline exactly why
theoretical predictions are necessary for us to understand more clearly what exactly
occurs within high-energy collisions at the LHC. Both these projects contribute significantly
to the jet and QCD community while offering direct and in-direct ways of impacting experimental
physics.

In chapter~\ref{chapter:DIS}, we introduced a new state-of-the-art jet algorithm 
for deep-inelastic scattering, which is a family of $k_t$ algorithms, which is currently going
through pier review. We presented an overview of the popular jet algorithms
used within the literature as while doing our research we realised there is not such
a comprehensive overview to understand how jet algorithms for DIS have been adapted.
We outlined the structure of the algorithm as implemented as the
\href{https://repo.hepforge.org/source/fastjetsvn/browse/contrib/contribs/DISGenkt}{\texttt{DISGenkt}}
plugin to \textsc{FastJet} as part of \texttt{fjcontrib}, taking special care to consider crucial points like
IRC safety how to order jets. We presented phenomenological results of this algorithm and compared to the
more recently presented algorithm, Centauro to demonstrate the differences our algorithm has.
Here we looked at the clustering sequences of the family of algorithms and Centauro and studies a range of
observables that gave us further insight into how the clustering formation occurs. Moreover, we looked at
hadronisation effects and observables that are particularly insensitive to hadronisation which can offer 
possible additional tests for QCD. Overall, this piece of work was really to introduce the algorithm given in 
\texttt{fjcontrib} and give some phenomenological results and compare to Centauro to demonstrate how it works.
We believe that this algorithm offers an exciting avenue to re-analyse the old HERA data, to see if this
algorithm sheds any light on anything missed previously. It has also come around at an exciting time where
the Electron-Ion Collider (EIC) which is due to undergo construction soon \todo{work out when it is},
this algorithm could take centre stage for analysing data due to the clean nature of separating the
two hemispheres. Additional extensions to this project could be in the form of resummation of jet rates. 
The \textsc{ARES} method introduced in Ch.~\ref{ch:hhZjet} could be used for this, where as discussed
in Sec.~\ref{sec:LTS-val}, this would need only an additional correction $\delta\mathcal{F}_{\rm clust}$
and only one initial-state radiation leg. \todo{Work out why we would need resummation for this project?}
This leads us on nicely to the second research project presented in this thesis.

In chapter~\ref{ch:hhZjet}, we presented an extension to the \textsc{ARES} framework for hadronic collisions.
This is work directed towards a future paper, where in this thesis we completed
the crucial conceptual stages of extending ARES to hadronic collisions which has never been done before. 
Here we have presented the formalism and have completed a mixture of the direct and indirect checks to validate our results so far, which have all been successful.


ARES: \\
Introduced novel resummation technique for hadronic collisions.
Produced the most conceptually difficult steps in this procedure.
Have provided a set of direct and indirect validations for the one-jettiness
and work has been started for other observables which aren't so simple 
to calculate in ARES. The matching to fixed order is the next bit. \\ \\


